\documentclass[11pt]{article}
\usepackage[margin=25mm]{geometry}
\usepackage{amsmath,amssymb,amsthm,booktabs,hyperref}
\newtheorem{theorem}{Theorem}
\newtheorem{lemma}{Lemma}
\newcommand{\Rea}{\operatorname{Re}}
\newcommand{\Ima}{\operatorname{Im}}
\newcommand{\tr}{\operatorname{tr}}
\title{Division-free reconstruction certificates at cutoff five}
\author{Research note}\date{7 October 2026}
\begin{document}\maketitle
\noindent\textbf{Status.} This note proves a finite-box exclusion theorem
and supplies a complete, independently replayable example. It does
\emph{not} prove the global implication $M(H)<5\Rightarrow d_F(H)<1/30$.
The missing step remains an exhaustive cover outside a proved Fourier
neighbourhood. Throughout, $M(H)$ denotes the maximum of the 400 corner
invariants of a raw order-six complex Hadamard matrix.

\section*{The certificate in one paragraph}
Choose a corner attaining $0<R=M(H)<5$. Our reconstruction gives six
linear equations for six bounded complex symmetric quantities on each
side. Maximality also gives twelve positive semidefinite $3\times3$
matrices. Multiply the equations by complex constants and the matrices
by positive matrices, then add. For an actual completion the result is
nonnegative. A certificate consists of rational multipliers for which
the same expression is strictly negative throughout a seed box, for
every allowed value of the bounded quantities. Thus the box contains
no such completion. Keeping the reconstruction equations undivided
avoids the small denominators in direct reconstruction enclosures.

\section*{1. The reconstruction equations}
Write the seed and complementary Gram matrix as
\[
 E=\begin{pmatrix}1&1&1\\1&a&b\\1&c&d\end{pmatrix},\quad
 G=\begin{pmatrix}3&\bar s&\bar t\\s&3&\bar r\\t&r&3\end{pmatrix},
 \quad s=-1-a-b,\ t=-1-c-d,\ r=-1-c/a-d/b.
\]
Here $(a,b,c,d)=(e^{i\theta_1},\ldots,e^{i\theta_4})$.
Use the reconstruction notation
\begin{align*}
 X&=|s|^2,&Y&=|t|^2,& Z&=|r|^2,&p&=X+Y+Z,\\
 J&=2\Rea(sr\bar t),& W&=XY+XZ+YZ-XYZ,\\
 g&=27-3p+J,&\eta&=9-p+J,&
 \chi&=W-(p-9)(p-18)+(p-15)J.
\end{align*}
Then $R=4g^2\eta-\chi^2$. Set $q=\sqrt R$ and $h=\sigma q$,
where $\sigma=\pm1$. Define
\begin{align*}
 P_0&=rs^2+\bar r t^2+(3-Z)st,\\
 P_1&=t\bar s^2+r^2\bar t+(3-Y)r\bar s,\\
 P_2&=s\bar t^2+\bar r^2\bar s+(3-X)\bar r\bar t,\\
 (k_0,k_1,k_2)&=(9-Z,9-Y,9-X),\qquad A_j=\chi+2k_j\eta,\\
 F_{ij}(h)&=2g^2+2k_i k_j\eta+(k_i+k_j)\chi+i(k_j-k_i)h.
\end{align*}
For normalized side columns $v_l=(1,x_l,y_l)^T$, put
\[
 u=\prod_lx_l,\quad v=\prod_ly_l,\quad w=u/v,\quad
 \Pi=st-\sum_lx_ly_l,\quad C_1=u\bar\Pi,\quad C_2=v\bar\Pi.
\]
The reconstruction identities, after clearing denominators, are
\begin{align}
 2\bar P_0P_1u&=F_{01}(h),&2\bar P_0P_2v&=F_{02}(h),&
 2\bar P_2P_1w&=F_{21}(h),\label{eq:products}\\
 2\bar P_0\Pi&=A_0+ih,&2P_1C_1&=A_1-ih,&
 2P_2C_2&=A_2-ih.\label{eq:permanent}
\end{align}
All coefficients are Laurent polynomials in the four seed phases.
The equations are affine in the real and imaginary parts of the six
unknowns and in $q$.

\begin{lemma}
For a physical side, the six absolute values are bounded by
$(1,1,1,3,3,3)$.
\end{lemma}
\begin{proof}
The first three are one and $|\Pi|^2=\eta$. The two complementary
permanent identities give $p+|J|\le9$, while
$|J|\le2(p/3)^{3/2}$. If $p\le27/4$, the latter is at most $p$.
If $p\ge27/4$, use $|J|\le9-p\le p$. Thus $J\le p$,
$\eta\le9$, and $|C_1|=|C_2|=|\Pi|\le3$.
\end{proof}

\section*{2. Positive matrices supplied by maximality}
The complementary-corner square identity supplies a real signed
invariant $f(T)=i\det L(T)$ for three phase columns, with
$f(T)^2=R(T^{\rm comp})$ in any physical Hadamard. The precise
determinant and its Laurent expansion are those in
\texttt{ProvenResults/SignedCornerPSDPrune/PROOF.md}.
For every pair of known seed columns $j<k$, define
\[
 N_{jk}=\sum_l f(E_j,E_k,v_l)v_lv_l^\dagger.
\]
Maximality gives $|f(E_j,E_k,v_l)|\le q$. Therefore
\begin{equation}\label{eq:psd}
 qG\pm N_{jk}
 =\sum_l\bigl(q\pm f(E_j,E_k,v_l)\bigr)v_lv_l^\dagger\succeq0.
\end{equation}
These are six positive matrices on each side. Their entries are affine
in the six complex symmetric quantities. This follows by expanding
the signed invariant and using Newton identities; for example
\[
 \sum_lx_l^2=s^2-2u\bar s,\qquad
 \sum_lx_ly_l=st-\Pi,\qquad
 \sum_lx_l^2/y_l=s\bar r-C_1.
\]
The actual 31 moment formulas and all six determinant coefficients
used by the implementation are checked as exact Laurent identities
by \texttt{verify\_dual\_algebra.py}.

Apply the same construction to $E^T$. The fixed-corner exhaustion and
matching theorem says that, outside Tao, its side-constructor sign is
$-\sigma$. This is a reconstruction input, not an inference from
positive-sheet surjectivity. The certificates below exclude Tao seeds
by an outward cube-root test on a seed coordinate. Hence only the two
matched sign pairs need checking.

\section*{3. A short certificate theorem}
Collect the twelve complex quantities as $\xi_1,\ldots,\xi_{12}$,
with radii $(r_1,\ldots,r_{12})=(1,1,1,3,3,3)$ repeated twice.
Let $\mathcal F_a\succeq0$ be the twelve matrices in
\eqref{eq:psd}, and $\mathcal E_b=0$ the twelve complex equations
\eqref{eq:products}--\eqref{eq:permanent}, on both sides.

\begin{theorem}[Finite-box exclusion]
Fix a seed box and a sign pairing licensed by the matching theorem.
Choose constant Hermitian positive semidefinite matrices $Y_a$ and
complex numbers $\lambda_b$. Expand the real affine expression
\[
 L=\sum_a\tr(Y_a\mathcal F_a)
       -\Rea\sum_b\bar\lambda_b\mathcal E_b
   =c_0+\sum_{j=1}^{12}
      (c_{jr}\Rea\xi_j+c_{ji}\Ima\xi_j)+c_q q.
\]
If throughout the box
\begin{equation}\label{eq:support}
 S(c):=c_0+\sum_{j=1}^{12}r_j\sqrt{c_{jr}^2+c_{ji}^2}
                 +\sqrt5\max(c_q,0)<0,
\end{equation}
then no physical completion in that box has that sign pairing and
$0<R(E)=M(H)<5$.
\end{theorem}
\begin{proof}
For such a completion, each trace in $L$ is nonnegative and each
equation term vanishes. Thus $L\ge0$. On the other hand, the radius
bounds and $0\le q\le\sqrt5$ imply $L\le S(c)<0$, a contradiction.
\end{proof}

The two sign certificates exclude all completions allowed by matching.
No side roots, side permutations or individual completion labels
need be enumerated. Dropping relations such as $w=u/v$ when taking
the upper bound enlarges the allowed domain and is therefore safe.

\paragraph{Retaining the radical relation.}
The exact relation $q^2=R$ supplies two further necessary conditions:
\[
 \begin{pmatrix}R&q\\q&1\end{pmatrix}
 =\begin{pmatrix}q\\1\end{pmatrix}(q,1)\succeq0,
 \qquad \frac94q-R=q\left(\frac94-q\right)\ge0.
\]
Here $q\le\sqrt5<9/4$. They are affine in the same unknown $q$;
their other coefficients are seed polynomials. Thus a positive
$2\times2$ matrix multiplier and a nonnegative scalar multiplier
can be added to $L$, without changing the certificate theorem.
The checker embeds the $2\times2$ matrix in a $3\times3$ matrix with
zero last row and column. It rounds and checks its multiplier exactly
as for the other positive matrices.

Nonnegative multiples of the canonical-order differences and
$9-p\pm J$ can be added too. Finally, the complementary signed
invariant satisfies $q\pm f(E)\ge0$ by maximality. Each such extra
term changes only the constant coefficient and, where applicable,
the coefficient of $q$. All its multipliers are nonnegative dyadic
rationals. These are necessary constraints only: the enlarged system
is not asserted to characterize actual Hadamard completions.

\section*{4. Checking an entire box}
The optimizer proposes $Y,\lambda$ only. The checker rounds their
entries to Gaussian dyadic rationals with denominator $2^{40}$,
makes $Y$ exactly Hermitian, and adds a small positive dyadic multiple
of the identity. Its leading principal minors are verified positive
using outward arithmetic. All subsequent polynomial coefficients
are exact Gaussian integers divided by powers of two.

For a real coefficient polynomial
$P(\theta)=\sum_\nu a_\nu e^{i\nu\cdot\theta}$, the degree-$d$
Taylor remainder on a cube of half-width $t$ is bounded by
\[
 E_P=\frac{t^{d+1}}{(d+1)!}
       \sum_\nu (|\Rea a_\nu|+|\Ima a_\nu|)\|\nu\|_1^{d+1}.
\]
This is the ordinary Taylor bound along the line from the centre
to the evaluation point. All centre derivatives are evaluated with
outward trigonometric arithmetic. Write
\[
 B(c)=|c_0|+\sum_jr_j\sqrt{c_{jr}^2+c_{ji}^2}+\sqrt5|c_q|.
\]
Subadditivity gives an upper bound for $S(c)$ by the centre support,
the sum of $t^{|\alpha|}B(\partial^\alpha c)/\alpha!$ for the
nonconstant Taylor terms, and $B(E)$ for the remainder.

A sharper bound keeps the quadratic part together. Write that part
in the tensor Bernstein basis of degree two on $[-t,t]^4$. The
one-dimensional coefficient vectors of $1,x,x^2$ are respectively
\[
 (1,1,1),\qquad(-t,0,t),\qquad(t^2,-t^2,t^2).
\]
The Bernstein basis functions are nonnegative and sum to one.
Consequently the quadratic coefficient vector lies in the convex hull
of its 81 Bernstein coefficient vectors. Since $S$ is convex, its
support is at most the largest $S$ of those vectors. Add the bounds
for Taylor degrees three through five and the sixth-order remainder.
This preserves the common phase displacement across the reconstruction
quantities and is substantially sharper than separate derivative bounds.

The same argument applies to the complete cubic Taylor polynomial,
using $4^4=256$ Bernstein control vectors. The higher-order checker
\texttt{higher\_support.py} uses this hull and adds only the fourth-
and fifth-order term bounds and the sixth-order remainder. Its
Bernstein conversion coefficients are generated and cross-checked
as exact rational numbers.

\section*{5. A complete retained example}
Consider the seed cube with centre and half-width
\[
 \theta_*={\pi\over64}(41,-33,-45,41),\qquad t={\pi\over64}.
\]
It is the depth-24 leaf with binary path
\texttt{100110010110110101000110} of the global seed partition.
Bisect each coordinate once, giving sixteen children of half-width
$\pi/128$. Four miss the necessary canonical Gram-modulus ordering.
The other twelve have both sign pairings excluded by the theorem.
Every verified support upper bound is less than $-0.49$.
Thus this whole canonical seed box contains no strict maximal carrier
with $M<5$. This is a box theorem, not a statement about sampled centres.

The complete numerical proposals are stored in
\texttt{dual\_parent\_complete.json}. They can be checked without
running an optimizer:
\begin{quote}\ttfamily\small
python3 simplerproof/averaging\_lemma/dual\_parent\_cover.py\par
\quad --reuse dual\_parent\_complete.json --replay-only\par
\quad --output dual\_parent\_replay.json
\end{quote}
The verifier regenerates each dyadic box, every polynomial, positivity
test and Taylor bound. It checks that all sixteen children are present.

\section*{6. What remains for the desired proof}
The local theorem in \texttt{PRACTICAL\_LOCAL.tex} already proves
\[
 \|u\|_F\le0.3,\quad H\sim F_6\circ e^{iu},\quad
 u\text{ horizontal},\quad M(H)<5
 \ \Longrightarrow\ d_F(H)<1/30.
\]
Subsequent full-domain Bernstein pruning, documented in
\texttt{GLOBAL\_COVER.tex}, reduces the open partition to 61,396 boxes.
The radical-enhanced certificates additionally exclude four full-size
parents, using the cubic support bound. Their worst support upper
bounds across the two signs are
\begin{center}\small
\begin{tabular}{lr}\toprule
Parent path & Certified upper bound, rounded upward\\\midrule
\texttt{111110010001110100010001}&$-0.449$\\
\texttt{110110110011101101111111}&$-1.054$\\
\texttt{110100111001010001100100}&$-12.880$\\
\texttt{110100110001110011101101}&$-11.258$\\\bottomrule
\end{tabular}
\end{center}
They are independently replayed by \texttt{verify\_radical\_dual.py},
which runs no optimizer. Their complete proposals, exact source hashes
and outward results are bound by \texttt{dual\_frontier\_checks.json}.
Two other tested parents remain open; their failed tests receive no
coverage credit.

Thus \textbf{61,392 boxes remain open}. A complete proof must exclude
all remaining remote boxes and prove that every other surviving
completion enters the local theorem's phase ball. The certificate
theorem supplies an exclusion tool, not a substitute for either obligation.

\paragraph{Dependencies.}
The reconstruction identities, side exhaustion and fixed-corner
matching are imported from the reconstruction proof dossier. The
complementary signed-square identity is imported from the corner-bound
proof. Their status and dependencies are distinct from the new exact
moment checks and finite-box certificates. Nothing here claims an
independent proof of those global reconstruction theorems.
\end{document}
